\chapter*{Agradecimentos e Informações Complementares}
\thispagestyle{empty}

Antes de tudo, deixo meu mais profundo agradecimento à minha família, minha irmã Paloma e especialmente ao meu pai, Paulo, e à minha mãe, Janice. Foram eles que me ensinaram, desde cedo, o valor do estudo e do esforço. Com amor, paciência e apoio incondicional, me deram a base necessária para que eu pudesse correr atrás dos meus objetivos. Tudo o que conquistei até aqui carrega um pouco de vocês.

Aos meus amigos, em especial àqueles que caminharam comigo durante toda a faculdade, Bárbara, Gabriel, Luan, Lucas, Marcus, Mateus e Matheus. Entre madrugadas estudando para provas difíceis e risadas nos intervalos das aulas, construímos muito mais do que trabalhos e memórias acadêmicas: construímos amizades que levarei comigo para a vida toda. Obrigado por tornarem essa jornada mais leve e mais divertida.

A todos os professores, servidores e funcionários terceirizados que, de alguma forma, fizeram parte da minha formação, deixo meu muito obrigado. Cada aula, cada orientação, cada gesto de apoio contribuiu para que eu crescesse não só como estudante, mas também como pessoa. Esse trabalho é, de certa forma, resultado do esforço coletivo de todos que cruzaram o meu caminho nesse período.

Ao meu orientador, Diego, agradeço pela confiança, pela paciência e por ter aceitado desenvolver este projeto ao meu lado, guiando cada etapa com dedicação e responsabilidade. À minha coorientadora, Thabatta, agradeço pelo acompanhamento, pelas contribuições e pelo apoio ao longo do desenvolvimento deste trabalho. Sem a orientação de vocês, esse projeto não teria tomado a forma que tomou.

Por fim, agradeço ao CEFET-MG, instituição que me acolheu desde 2019, quando ingressei no curso técnico de Informática para Internet e passei a enxergar a tecnologia não apenas como estudo, mas como vocação. Foi ali que me apaixonei pela área, pelo campus e pela vida acadêmica, e onde vivi experiências que marcaram profundamente a minha trajetória. Levo comigo não só a formação, mas também as memórias e aprendizados construídos ao longo desse caminho.

\vspace{0.5cm}
\noindent\textbf{Transparência no uso de ferramentas de IA}

O autor informa que contou com o auxílio das ferramentas ChatGPT (versão 5.5) e Gemini (versão 3.1) como apoio na revisão gramatical, no aprimoramento da fluidez do texto e na sugestão de estrutura, bem como no suporte à elaboração de partes do código e à definição de casos de teste e verificações experimentais. Todo o conteúdo produzido com esse apoio foi revisado, validado e adaptado pelo autor, sendo a responsabilidade final pelo trabalho exclusivamente do autor.